\chapter{绪论}\label{chap:intro}
\section{研究背景与问题提出}
异常检测关注不符合预期行为的观测或模式。相关研究在统计假设、异常定义和
应用约束上存在差异，不能仅凭一个分数阈值统一解释所有场景\cite{chandola2009survey}。
本例将无人机飞行记录抽象为多变量时间序列：每个架次由同步采集的多个变量组成，
不同架次之间则保留独立的起止边界。这样的表示便于讨论时序依赖、变量关联与部署流程。

对于同一时刻的观测，高残差既可能反映异常，也可能来自模型未学习到的正常机动。
因此，论文中的问题描述应区分观测空间、模型假设和决策规则。
本例关注的问题不是“用复杂模型得到更高分数”，而是怎样把数据划分、预测、
阈值校准和指标计算组成一条清晰可复查的链路。

\subsection{示例目标与边界}
本文具有两个相互区分的目标。第一，给出可执行的小型时序异常检测流程；
第二，提供覆盖学位论文常见元素的排版样本。两者通过统一的数据文件和自动生成的
结果表连接，但不能相互替代。模板编译成功并不意味着模型有效，
合成数据上得到有限的检测效果也不意味着模板满足所有学校规范。

示例没有调用真实无人机数据集，也没有实现神经网络预训练。
“飞行”“传感器”“故障”等词在实验部分描述的是合成信号的设计角色。
所给出的方法为透明基线，最终结论限于该脚本和指定参数。
这一约束也适用于封面、摘要与英文摘要中的表述。

\subsection{工程组织原则}
论文源码按章节拆分，实验数据与图表由脚本生成，参考文献统一维护在一个数据库中。
正文不手写“图 3-1”或“式（4-2）”等会随编辑失效的编号，而使用标签和交叉引用。
所有浮动体的题注先于标签写入，使引用指向相应对象，而不是此前的章节计数器。

\section{相关研究概述}
\subsection{序列预测与表示学习}
长短期记忆网络通过门控结构处理序列中的信息传递问题\cite{hochreiter1997lstm}。
Transformer 则以注意力机制组织序列表示\cite{vaswani2017attention}。
这些方法提供了不同的建模工具，但模型结构与异常判据不是同一概念。
将预测器用于异常检测时，仍需定义残差聚合、训练数据范围及阈值来源。

\citet{nie2023patchtst} 通过分块输入和通道独立设计研究时间序列预测。
\citet{zeng2023dlinear} 则以简单线性模型检验一组预测基准上的结构选择。
这些研究说明，比较对象应与具体任务和实验协议对应；
本例选择岭回归只是为了降低复现门槛，不是用未经运行的神经网络结果充当对照。

\subsection{异常判据与通用时序模型}
Deep SVDD 研究在特征空间中约束正常样本表示的方法\cite{ruff2018deep}。
Anomaly Transformer 通过关联差异构造时序异常判据\cite{xu2022anomaly}。
MOMENT 将多个时序任务纳入预训练模型框架\cite{goswami2024moment}。
本例在综述中引用这些工作，但实验中不报告它们的性能，
也不将自己的合成数据分数与其论文中的公开结果直接排序。

\subsection{软件工具与复现}
机器学习基础概念可参见周志华的教材\cite{zhou2016ml}，
深度学习中的表示、优化与正则化讨论可参见 Goodfellow 等人的专著\cite{goodfellow2016deep}。
在工具层面，scikit-learn 提供了多种机器学习算法接口\cite{pedregosa2011sklearn}，
PyOD 将不同异常检测方法组织在统一工具箱中\cite{zhao2019pyod}。
本工程的数值脚本只依赖 NumPy 与 Matplotlib，不要求读者安装上述完整工具箱。

\section{本文内容与章节安排}
第~\ref{chap:foundations}~章定义时序窗口、异常分数及评价指标；
第~\ref{chap:data}~章说明合成数据和架次级划分；
第~\ref{chap:method}~章给出岭回归、平滑和阈值校准；
第~\ref{chap:experiments}~章报告自动生成的结果并讨论局限；
第~\ref{chap:conclusions}~章总结案例与后续工作。
附录~\ref{app:config}、附录~\ref{app:proof} 和附录~\ref{app:reproduce}
分别提供配置字典、补充推导与复现清单。

本节同时展示三种引用写法：普通引用\cite{chandola2009survey}、
行内引用\ExampleInlineCite{vaswani2017attention}，以及
\citet{goswami2024moment} 形式的叙述引用。
参考文献后端的切换只影响著录和引用格式，不改变章节内容。
